%% ==================================================================
%% From the research notes "Not Every Cut Is an Aperture" and
%% "Where Does Agency Live?"  --- the mechanism.
%% ==================================================================
\chapter{Not Every Cut Is an Aperture}
\label{ch:apr}

\noindent
The previous chapter built a tower and indexed it by ordinals, which is how anyone
first gets the intuition and is not how the thing actually works. Turing's
construction relativises a machine to an oracle, takes the halting problem of the
relativised machine as the next oracle, and iterates; the index is an ordinal because
the construction is a sequence. But an ordinal is a chain, and there is no reason the
capacities that separate one computational system from another should be linearly
ordered. This chapter and the next replace the chain with what is underneath it, and
the replacement begins by asking a question the tower never had to face: where in a
computation does the extra power actually get spent?

The answer proposed here is that it is spent resolving races, and that this is not a
metaphor. When two sends contend for one receive, something must pick. That picking is
a choice function for a definite family, the assertion that a complete schedule exists
is an instance of the axiom of choice, and the strength of the choice principle
required is graded --- free when the contention is finite, countable or dependent
choice at one $\omega$, full choice only for class-sized families. So the oracle was
never something to bolt on. In any calculus where interaction is primitive and
confluence fails, the oracle is already present at every cut, wearing the name
``scheduler.''

What follows is the instrument. It is deliberately mechanical: what a cut is, why the
nondeterminism belongs to the cut rather than to either of the processes flanking it,
what an interior has to satisfy to carry no observable choice at all, and where the
boundary of such an interior falls. The next chapter turns the instrument on the
question of where an agent's boundary goes. Read in that order the machinery arrives
before the thesis it supports, which is the right way round even though the thesis is
the more interesting half.

\bigskip
\section{Interaction, and why non-confluence is the point}\label{ndo:sec:interaction}

Among the objects of $\mathbf{GSLT}$, one class is built for our question. The mobile
process calculi (MPCs) --- Milner's $\pi$-calculus and, in the form we use, the
reflective higher-order (rho) calculus \cite{RHO} --- take a definite philosophical
stance: \emph{all computation arises from interaction among agents}
\cite{milnerCC}. There is no lone machine grinding a tape; there are processes that
meet, exchange, and continue, and computation is what happens at their meetings. We
write the meeting as a parallel composition, the \emph{cut}:
\[
P \Par Q,
\]
and the single base reaction is communication across it,
\[
\inn{y}{x}{P} \Par \out{z}{Q} \reduces P\{@Q/y\} \qquad (x \equiv_N z),
\]
an input on a channel meeting a matching output. Agents flank every cut; the cut is
where they act on one another.

The feature that makes the MPCs the right tool --- and that distinguishes them
sharply from the $\lambda$-calculus --- is that they are \emph{fundamentally
non-confluent}. When several outputs and several inputs collect on the same channel,
which output meets which input is not determined by the terms; it is a real choice,
and different resolutions lead to genuinely different, non-reconvergent futures.
This is usually presented as a complication. We present it as the entire point,
because non-confluence is the computational signature of a \emph{race}, and races
are how agency shows up in the world. Consider, across four scales:

\begin{itemize}[leftmargin=1.4em]
\item two sodium atoms competing to donate their electron to one chlorine;
\item two sperm racing to fertilise one egg;
\item two predators converging on a single prey;
\item two buyers reaching for the last concert ticket.
\end{itemize}

\noindent
Each is, structurally, one consume-slot (the chlorine, the egg, the prey, the
ticket) contended by several produce-candidates, resolved by a pairing that the
situation does not fix, with the unmatched candidates left over. Each outcome
\emph{matters}: this sodium and not that one is now bound; this sperm and not that
one founds the lineage. A confluent model cannot speak of these, not because it
lacks the vocabulary but because confluence is precisely the claim that the choice
does not matter --- that all resolutions reconverge. The $\lambda$-calculus, the
great confluent model, is the stakes-free corner of the landscape: its many
reduction orders are real choices whose outcomes are identified by Church--Rosser,
so the choice is free and invisible. Reality is mostly not like that. To talk about
agency one needs a model in which the choice can matter, and the MPCs are the
objects of $\mathbf{GSLT}$ in which it does. This is the first crossing of our two axes:
biology (the race for the egg) tells us which computational models are adequate to
agency, and the answer is the non-confluent ones.

\paragraph{The rho calculus, specifically.}
We use the rho calculus because its names are quoted processes --- $@P$ turns a
process into a name, $\deref x$ turns a name back into a process, and
$\deref @P \equiv P$ --- so there are no primitive names and no name-restriction
operator; reflection furnishes recursion with no primitive replication \cite{RHO}.
Nothing below turns on the set-theoretic models of \cite{qcs,knot}; we draw only on
the shape of the calculus. The reader may keep the modern surface syntax in view:
$@P$ for quote, $\deref x$ for dereference, $\inn{y}{x}{P}$ for input, $\out{x}{Q}$
for output, $\Par$ for the cut, $\rcong$ for structural congruence, $\bisim$ for
context bisimilarity (the behavioural, congruence-by-construction equivalence) and
$\rbisim$ for resource bisimilarity (the finer equivalence that counts parallel
copies).


\section{RSpace: the bridge, and where the choice is}\label{ndo:sec:rspace}

The MPCs are admirable for reasoning and awkward to run. RSpace is the abstraction
that closes the gap: it presents the parallel fragment of the rho calculus as a
store --- a hashtable from channels to polarity-homogeneous bags, either of data
(produces) or of continuations (consumes) --- and interprets the cut as keyed
multiset union, so the denotation $\sem{-}$ is a monoid homomorphism and the
structural laws of $\Par$ become literal equalities of tables \cite{qcs}. This is
not a toy: it is how the calculus is compiled into a high-performance, persistent,
content-addressed store in production. We invoke it here for one reason: it makes
the choice of \S\ref{ndo:sec:interaction} \emph{visible and operational}.

When a channel carries both a bag of consumes and a bag of produces, communication
chooses an underlying list of each and \emph{zips} them --- matching pairs react,
unmatched payloads remain as surplus at the channel. The denotation stays a
function; the zip is the separate dynamics; and the choice of which list orderings
to zip is the calculus's nondeterminism, now pinned to a key. The four races of
\S\ref{ndo:sec:interaction} are this forgiving zip with surplus, drawn from life:
several candidates, one slot, a pairing the data does not fix, leftovers remaining.
RSpace's contribution to the present inquiry is to show that the nondeterminism is
not a haze over the system but a specific, locatable act of selection --- and an act
of selection is a choice, in the sense the next section makes exact.

\begin{remark}[A motivating example, not a mechanism]
We are not proposing that sodium, sperm, or ticket-buyers \emph{run} RSpace, or that
its store is a physical mechanism. The claim is structural: the abstract shape ---
contended slot, nondeterministic pairing, surplus --- is shared, and RSpace is where
that shape, and the choice inside it, is displayed cleanly enough to analyse. The
value of the MPC and RSpace vocabulary is that it \emph{names the categories};
whether any given physical system instantiates them is a separate, empirical
question we do not prejudge.
\end{remark}


\section{Choice as a computational resource}\label{ndo:sec:choice}

Begin with the one feature of the RSpace example we will actually use. When a
channel carries both a bag of consumes $\{\!|(y_1)P_1,\dots,(y_n)P_n|\!\}$ and a
bag of produces $\{\!|Q_1,\dots,Q_m|\!\}$, communication chooses an underlying
list of each and zips them, keeping unmatched payloads as surplus. The choice of
list orderings is the nondeterminism of the calculus --- which send meets which
receive --- now localised at the key. The denotation itself remains a function;
reaction is the separate dynamics; and the indeterminacy is exactly the choice of
pairing.

That is not a metaphor for choice. It \emph{is} a choice function, for a specific
family. Index by the cuts that carry mixed polarity (or, for a run, by the
reaction events); at each such index the fibre is the set of admissible pairings,
inhabited precisely when a redex is present. A completed schedule is a section of
this family --- an element of the product of the fibres --- and the assertion that
such a section exists is an instance of the axiom of choice. The scheduler is the
choice function for the pairing family, not by analogy but by construction. This
gives the practical reframing we want: \emph{choice is a computational resource},
the resource of nondeterminism, and different ways of supplying it import
different amounts of computational power.

\paragraph{The resource is graded by the same $\omega$ the risk ledger charges.}
The strength of the choice principle one must invoke is not uniform; it tracks
the infinitude of the family.

\begin{itemize}[leftmargin=1.4em]
\item \emph{Finite bags, finite run.} The index set is finite, the product is
inhabited with no choice axiom at all (constructively, in $\mathrm{IZF}$). More
than that: a bag is a parallel composition of finitely many prefixes, so it
arrives \emph{presented} as a list-up-to-permutation. The term hands you the
enumeration; picking a linearisation is constructively free. The term is its own
choice structure.

\item \emph{One $\omega$.} Replication and reflection install an $\omega$-fold
bag, or a non-terminating run yields an $\omega$-sequence of events. Now one needs
countable choice for the static store, or --- for a run whose pairings depend on
the residuals of the previous step --- \emph{Dependent Choice}. The cascade of a
deep cut unpacking into a deeper cut \cite{paths} is exactly a dependent-choice
sequence. This is the single $\omega$ --- the ``black infinity'' --- the risk
ledger of \cite{qcs} already spends; choice-as-resource and the $\omega$-meter
are the same gauge.

\item \emph{Class-sized.} Full $\mathrm{AC}$, and only here.
\end{itemize}

\noindent
So the slogan ``choice is not non-constructive'' is exactly true in the finite
fragment and \emph{relocated, not dissolved}, above it. The constructive escape
at $\omega$ is to never form the completed schedule: take each pairing locally, on
demand, as a decidable step, replacing the choice \emph{function} (whose bare
existence needs $\mathrm{DC}$ or $\mathrm{AC}$) with a choice \emph{process}. That
is what lazy, on-demand reaction does operationally; it keeps choice potential
rather than actual. The honest reading is of a piece with the rest of the
programme: recursion relocated from anti-foundation into reflection, congruence
from an afterthought into the context labels, and now choice from a metatheoretic
existence axiom into an operationally consumed resource.

\paragraph{A lattice, not a tower.}
The natural home for ``choice as a resource'' already exists: the Weihrauch
lattice of Brattka, Gherardi and Pauly \cite{weihrauch}, whose objects are
choice principles ordered by uniform reducibility, with demonic binary choice
appearing as the lesser-limited-principle-of-omniscience and composition
$\star$ modelling ``use one resource, then another.'' The decisive structural
fact, for our purposes, is that the Weihrauch degrees form a \emph{lattice} with
rich algebra, not a linear hierarchy. We will return to this: the intuition that
a universe of computational agents is ``of all stripes, not neatly sorted by a
tower of oracles'' is, formally, the statement that the relevant degree structure
is a lattice and not an ordinal chain.


\section{Apertures, and why the cut owns them}\label{ndo:sec:aperture}

Two notions first, since everything after this turns on them.

\begin{definition}[Cut and aperture]\label{ndo:def:aperture}
A \emph{cut} is a parallel composition co-locating a consume and a produce --- the
locus where two agents interact. A cut is an \emph{aperture} when the resolution of
its pairing nondeterminism is \emph{observable}: distinct admissible pairings lead to
residuals in distinct behavioural classes (distinct $\rbisim$-classes at the resource
layer, or distinct $\bisim$-classes after forgetting multiplicity). A cut that is not
an aperture carries no observable choice: either there is a unique match, or the
several matches reconverge.
\end{definition}

\begin{principle}[Not every cut is an aperture]\label{ndo:prin:main}
Agents flank every cut, but only apertures leave the interaction's nondeterminism
open. And only at an aperture can an agent of higher computational capacity influence
the behaviour of a lower one --- by biasing the live choice. Where a cut is not an
aperture there is no choice to bias, and the stronger agent's surplus power finds no
purchase. Agency, in the sense of room-to-do-otherwise, lives at the apertures.
\end{principle}


It is tempting to summarise \S\ref{ndo:sec:choice} by saying that an agent's
\emph{nondeterminism budget} is its interface to the rest of the universe: a
fully deterministic process is sealed, with every slot filled by its own
transition function and nowhere for anything external to write; a
nondeterministic agent is open precisely at its choice points. There is a true
idea here --- under-determination is the aperture --- but the locus is wrong, and
the parallel-composition algebra is what makes it wrong.

Nondeterminism is not preserved by $\Par$ in either direction. A whole that is
deterministic can decompose into parts each wildly underdetermined at its
\emph{internal} cuts; and parts each internally rigid can compose into a whole
whose only indeterminacy is at the boundary between them. The aperture is
therefore a property of the \emph{cut}, not of either agent at it, and --- this is
the consequence that matters --- \emph{which} cut one privileges as ``the agent
boundary'' is a modelling choice. Different observers draw it at different depths
of the same term, and the interior lights up or goes dark accordingly.

\begin{remark}[Under-determination versus boundary]\label{ndo:rmk:seam}
Apertures coincide with under-determination: a cut is an aperture exactly when the
resolution of its pairing nondeterminism is \emph{observable}, i.e.\ when distinct
admissible pairings yield residuals in distinct $\rbisim$-classes (equivalently,
distinct $\bisim$-classes, after the projection that forgets multiplicity). But
the seam at which under-determination sits need not respect any particular
delineation of agents. A seam may lie far above, or far below, the boundary an
observer would draw; its effects may be felt at cuts the observer counts as
internal. This is why ``the agent's budget is its interface'' fails: the budget is
real, but it belongs to the cut, and the cut may straddle whatever boundary one
had in mind.
\end{remark}

In the path reading of \cite{paths}, where keys are routes through a trie and the
prefix order gives one channel a subspace of another, the modelling freedom
becomes literal: \emph{a choice of agent boundary is a choice of prefix}, and the
trie is the lattice of possible boundaries. A coarse boundary is a short prefix,
near the root, folding everything below into one sub-store and exposing only a
shallow aperture; a fine boundary is a long prefix, exposing apertures at every
internal branch. These are not rival claims about where the boundary ``really''
is. They are different truncation depths on one structure, ordered by the prefix
order, which is precisely the order ``coarser agent / finer agent.'' The In-deep
rule of \cite{paths} --- a reaction at a shallow cut unpacking a cut that lives
deeper than the boundary-drawer was tracking --- is what it looks like when the
same dynamics is read at two depths at once.


\section{The macro-component: confluence, not linearity}\label{ndo:sec:macro}

We want a name for the object Smith's interior instantiates: a sub-term all of
whose internal cuts are sealed, with apertures only at its boundary. The first
guess is that the interior is sealed by \emph{rigidity} --- that at every internal
key the consume-bag and produce-bag admit a \emph{unique} compatible pairing, so
that the forgiving zip is a function and there is no race to resolve. Determinism
not by scheduling but by there being nothing to schedule. Syntactically this is a
\emph{linearity} condition: every internal channel is used at multiplicity one per
polarity, each consume a linear hypothesis discharged by exactly one produce, the
boundary the residue of the non-linear, contended channels. On this reading
\textsf{EnergyEaters} are exactly the non-linear channels --- the ones contended by
many consumers, the sun feeding everyone --- and the factorisation is
\emph{linear interior, non-linear boundary}.

Linearity is, however, too strong. It is a syntactic sufficient condition for what
we actually care about, which is that the interior have no \emph{observable} choice.
The semantic statement of exactly that, and nothing more, is confluence.

\begin{definition}[Rigid sub-term / macro-component]\label{ndo:def:macro}
A sub-term is \emph{rigid} (internally) if its internal reaction is confluent up to
resource bisimilarity: every internal pairing-race reconverges, so all maximal
internal reaction sequences reach $\rbisim$-equal normal forms. A
\emph{macro-component} is a maximal rigid sub-term: every internal cut is a
non-aperture (confluent up to $\rbisim$), and its boundary cuts --- where
confluence fails --- are its apertures. The set of cuts at which internal
confluence fails is the \emph{rigidity frontier}.
\end{definition}

\noindent
Confluence permits the branching that linearity forbids. The bee may have several
internally admissible pairings; if they all flow to $\rbisim$-equal normal forms,
the bee is still not a site of observable choice, which is all the factorisation
needs. Metabolism is full of internal branching --- parallel pathways, redundant
enzymes, order-independent cascades --- that converges on the same steady state:
confluent, not linear. So the right slogan is \emph{the interior is
$\lambda$-like, not linear-like}.

\paragraph{$\lambda$ is really deterministic.}
The $\lambda$-calculus is the load-bearing witness, not an aside. Operationally it
is radically nondeterministic: a term with many redexes has many reduction
sequences, every interleaving a distinct branch. And it is \emph{the} paradigm of
determinism --- because Church--Rosser says the branching is fictional: one normal
form, reached however you like. The choice of redex is real and it is free in the
sense of \S\ref{ndo:sec:choice}, semantically free, spent but unobservable, eliminated
by the quotient to the normal form. $\beta$-reduction is the cleanest instance in
all of computing of a process saturated with races whose races do not matter, and
it is deterministic for precisely the reason Definition~\ref{ndo:def:macro} asks of an
interior: confluence, not absence of branching. A \emph{closed} $\lambda$-term ---
no free variables --- is then a sealed macro-component at every scale, nothing to
interfere with because there is no contended channel for anything to write to. An
\emph{open} term is the agent with an aperture: its free variables are its boundary
channels, the spots where a context (a higher agent) can substitute and change the
result. In the MPC vocabulary the free variable is the unmatched consume, the
dangling polarity --- the same object, the port where the universe reaches in.

\begin{proposition}[Degree is carried by the boundary]\label{ndo:prop:degree}
A confluent (up to $\rbisim$) interior contributes no choice resource: it carries
no Weihrauch content, no observable bits of nondeterminism, however baroque its
branching, because the resource was defined (\S\ref{ndo:sec:choice}) as observable
choice and confluence makes choice unobservable. Hence, modulo the obligations
below, the resource degree of a macro-component is carried entirely by its
boundary:
\[
\deg(\text{macro-component}) \;=\; \deg(\text{boundary}),
\]
the interior being degree-transparent. A linear-interior criterion would
under-count here, miscalling some genuinely confluent interiors ``boundary''
because they branch; the confluence criterion is what makes the equation clean.
\end{proposition}


\subsection{The islands are small}

The picture this forces is not one determinate interior behind one skin. If
selection, predation, and competition are non-confluent and pervasive --- and
\S\ref{ndo:sec:smith} will grant that they are --- then confluent macro-components are
\emph{small and local}: a tightly coupled metabolic or developmental subsystem whose
internal branching genuinely reconverges, redundant enzymatic pathways and buffered
developmental cascades, bounded by an aperture wherever a real race begins. The
biosphere is then a vast field of small confluent islands in a dense sea of races,
rather than a single interior. This is less tidy than a clean interface and it is the
true shape: \emph{apertures are everywhere}, the determinate regions are the exception,
and ``where is the agent?'' has many local answers rather than one global one. An
ecosystem may be the right scale for some such island --- but only the parts of its food
web that reconverge up to resource bisimilarity, with its contended predation and its
competition for influx counted as apertures \emph{at or within} it, not as interior. The
ecological races are not quarantined behind the boundary; they are the boundary,
threaded through.

Here is where two senses of ``agent'' come apart, and the distinction is the conceptual
payoff.

\begin{remark}[Two senses of ``agent'']\label{ndo:rmk:two}
A confluent island's internal parts may be sharply individuated --- distinct enzymes,
distinct cell types, distinct roles --- without being \emph{sites of choice}. Their
distinctness is type-distinctness, not aperture-distinctness; whatever internal
branching they run reconverges, so no internal cut is a live race. Discrimination among
agents and the location of nondeterminism therefore come apart: an agent in the
\emph{discrimination} sense, an individuated role, need not be an agent in the
\emph{choice} sense, a cut with a live race --- and conversely a featureless region can
sit astride a fierce one. A finer description sees \emph{more} individuated agents but
not necessarily more apertures; the apertures are wherever the races are, which is a
separate question from where the roles are. Naive ontology conflates the two, placing
``the agent'' at the individuated role.
\end{remark}

\begin{remark}[The same boundary, priced]\label{ndo:rmk:ecology}
The chapters on ecology reached a criterion for individuation from the opposite
direction, and it is worth putting the two side by side because they disagree in an
informative way. There, an individual is a region that may close its namespace ---
partially, since \emph{perfect closure is death}
(\S\ref{eng:S5.5} of \chEng) --- and the criterion is economic: close where repair is
cheaper than exposure. Here the criterion is semantic: close where no internal cut
carries observable choice. Neither implies the other. A region can be metabolically
worth owning and full of live races, which is most organisms; a region can be
determinate and not worth closing, which is most enzymes. What the two share is the
refusal to put the boundary at the organism by default, and the observation that
whichever criterion is used, the boundary is porous --- an aperture and a channel with a
redex pathway through it are the same object described in two vocabularies.
\end{remark}

\section{Three frontiers, one locus}\label{ndo:sec:frontier}

The reward of taking confluence as the criterion is that the boundary of the
macro-component coincides with two frontiers we already had names for.

\begin{proposition}[Frontier coincidence]\label{ndo:prop:frontier}
The following three loci coincide, for a macro-component in the sense of
Definition~\ref{ndo:def:macro}:
\begin{enumerate}[label=(\roman*),leftmargin=2.2em]
\item the \emph{rigidity frontier} --- the cuts where internal confluence up to
$\rbisim$ fails;
\item the \emph{detectable-interference frontier} --- the cuts where a biased
pairing yields a different $\rbisim$-class, so that an external agent's write
leaves an observable mark rather than being absorbed into a $\rbisim$-equal
outcome;
\item the \emph{confluence obligation frontier} --- the cuts at which the
eager-confluence-up-to-$\rbisim$ obligation of the path semantics
(\cite{paths}, Obligation~2) is non-trivial.
\end{enumerate}
\end{proposition}

\noindent
The coincidence is a consistency check on the vocabulary: the place where the
interior stops being able to hide its schedule is the place where external
interference starts being able to leave a mark, which is the place where the
formal confluence obligation bites. One frontier, three readings ---
confluence-failure from inside, interference-visibility from outside, proof
obligation from the meta-theory. That a single locus answers to all three is the
kind of agreement that suggests the apparatus is carving at a joint.

It also tells us when interference is a \emph{measurement}. Where internal reaction
is confluent up to $\rbisim$, a biased pairing picks the trajectory but not the
destination; the weak agent is steered to a $\rbisim$-equal normal form and cannot
tell. Where confluence fails, the bias becomes a different observable outcome.
Detectable interference is interference at a non-confluent cut --- an interaction
that resolves an otherwise-unobservable choice into an observable distinction,
which is close to the definition of a measurement. We note, without leaning on it,
that the way a weak agent could in principle detect a stronger one is
information-theoretic and contextual: its ``free'' choices would look random
relative to its own capacity yet be compressible relative to the higher degree,
carrying more mutual information with the oracle than chance allows --- local
sections admitting no degree-uniform global section, in the
sheaf-theoretic sense of contextuality \cite{abramsky}. The probabilistic face is
cleaner: a flat jumble is one joint distribution, an agent's view a marginal, and a
strong agent biasing a weak one's choices is conditioning that marginal on oracle
information the weak agent cannot access. We flag these as directions, not results.


\section{Choice types, and the terms that inhabit an aperture}\label{ndo:sec:types}

Everything above treats choice as a quantity: how much of it a cut consumes, where in
the Weihrauch lattice that lands. It can be treated as a \emph{type} instead, and
\chChs\ does exactly that --- a Curry--Howard correspondence in which choice principles
are the types and fair schedulers are the terms. Two of its consequences are needed
here, ahead of their development.

The first is a reading of the macro-component. At an aperture the obligation is a
$\forall\exists$ statement: for the contended slot, \emph{there exists} an admissible
pairing. That formula is the \emph{choice type} of the aperture, and a scheduler
producing the pairing is a proof of it. A macro-component in the sense of
Definition~\ref{ndo:def:macro} is then a region whose choice type is
\emph{proof-irrelevant}: several pairings inhabit it, all normalise to $\rbisim$-equal
results, so the term carries no observable information. ``Capacity lives at the
apertures'' becomes ``the term content lives where the choice type is not
proof-irrelevant'' --- the same claim with a metatheory attached to it.

The second is why the next chapter's closing question is a question about typing rather
than about dynamics. \emph{Subject reduction} --- that running a schedule preserves the
type --- says operationally that scheduling cannot bootstrap an oracle. Whether it holds
is the load-bearing open problem of this part, and \S\ref{ndo:sec:crux} states it in
both dresses.
